\thispagestyle{plain}
\begin{center}
	\vspace*{1.5cm}
	{\Large \bfseries  Resumen}
\end{center}
\vspace{0.5cm}
Para optimizar la toma de decisiones en la compra de productos de la empresa, es fundamental desarrollar un modelo predictivo de demanda que integre datos históricos de ventas con información relevante, como precios, cantidades, días festivos, número de transacciones y otros factores. Este proceso comienza con una recopilación exhaustiva de datos, seguida de un análisis exploratorio para identificar patrones y tendencias significativas.  

La selección de características permite determinar las variables que tienen mayor impacto en la demanda. Posteriormente, se emplean técnicas de modelado predictivo, como regresión lineal, redes neuronales, máquinas de soporte regresion (SVR), Random Forest, ARIMA, Prophet, entre otros modelos de forecasting y aprendizaje automático, para predecir el volumen futuro de ventas. Una vez entrenado, el modelo se valida rigurosamente para garantizar su precisión.  
Tras su implementación, es crucial realizar un seguimiento continuo del modelo para ajustarlo y optimizar las predicciones frente a cambios en el mercado y otros factores relevantes. Esto permite mejorar la gestión de decisiones estratégicas y diseñar políticas de precios más efectivas.  
La investigación se centra en el uso de modelos de series de tiempo para identificar las áreas con mayor volumen de ventas, maximizando así los ingresos potenciales. Los modelos de series temporales son herramientas analíticas que permiten analizar y predecir patrones en datos secuenciales a lo largo del tiempo.  

Para llevar a cabo este enfoque, es esencial recopilar datos detallados de ventas por familia de productos, incluyendo información como volúmenes de ventas diarios, semanales o mensuales, ingresos generados y factores externos, como promociones, eventos especiales, fines de semana o días festivos. Mediante el análisis de estos datos, es posible identificar patrones y tendencias que faciliten pronósticos precisos para la empresa Donna S.A.C.  

El análisis de series de tiempo implica estudiar la estructura temporal de los datos, considerando factores como tendencia, estacionalidad y variaciones aleatorias. Con esta comprensión, se aplican modelos estadísticos y técnicas de pronóstico para estimar el comportamiento futuro de las ventas. Utilizando este enfoque, la empresa puede obtener estimaciones confiables y precisas que impulsen una mejor toma de decisiones estratégicas y operativas.\\

\textbf{Palabras claves: } Predicción, Ventas, Machine Learning, Regresión